

Now, the definition of the infinite product $\prod_{i\in I}\mathbf{a}_{i}$
(namely, Definition \ref{def.fps.multipliable} \textbf{(b)}) yields that%
\begin{equation}
\left[  x^{n}\right]  \left(  \prod_{i\in I}\mathbf{a}_{i}\right)  =\left[
x^{n}\right]  \left(  \prod_{i\in M}\mathbf{a}_{i}\right)
\label{pf.prop.fps.union-mulable.b.2}%
\end{equation}
(since $M$ is a finite subset of $I$ that determines the $x^{n}$-coefficient
in the product of $\left(  \mathbf{a}_{i}\right)  _{i\in I}$). On the other
hand, the set $U$ is an $x^{n}$-approximator for $\left(  \mathbf{a}%
_{i}\right)  _{i\in J}$. Thus, Proposition \ref{prop.fps.infprod-approx-xneq}
\textbf{(b)} (applied to $J$ and $U$ instead of $I$ and $M$) yields
\begin{equation}
\prod_{i\in J}\mathbf{a}_{i}\overset{x^{n}}{\equiv}\prod_{i\in U}%
\mathbf{a}_{i} \label{pf.prop.fps.union-mulable.b.3}%
\end{equation}
(since the family $\left(  \mathbf{a}_{i}\right)  _{i\in J}$ is multipliable).
Furthermore, the set $V$ is an $x^{n}$-approximator for $\left(
\mathbf{a}_{i}\right)  _{i\in I\setminus J}$. Thus, Proposition
\ref{prop.fps.infprod-approx-xneq} \textbf{(b)} (applied to $I\setminus J$ and
$V$ instead of $I$ and $M$) yields
\begin{equation}
\prod_{i\in I\setminus J}\mathbf{a}_{i}\overset{x^{n}}{\equiv}\prod_{i\in
V}\mathbf{a}_{i} \label{pf.prop.fps.union-mulable.b.4}%
\end{equation}
(since the family $\left(  \mathbf{a}_{i}\right)  _{i\in I\setminus J}$ is
multipliable). From (\ref{pf.prop.fps.union-mulable.b.3}) and
(\ref{pf.prop.fps.union-mulable.b.4}), we obtain%
\[
\left(  \prod_{i\in J}\mathbf{a}_{i}\right)  \left(  \prod_{i\in I\setminus
J}\mathbf{a}_{i}\right)  \overset{x^{n}}{\equiv}\left(  \prod_{i\in
U}\mathbf{a}_{i}\right)  \left(  \prod_{i\in V}\mathbf{a}_{i}\right)
\]
(by (\ref{eq.thm.fps.xneq.props.b.*}), applied to $a=\prod_{i\in J}%
\mathbf{a}_{i}$ and $b=\prod_{i\in U}\mathbf{a}_{i}$ and $c=\prod_{i\in
I\setminus J}\mathbf{a}_{i}$ and $d=\prod_{i\in V}\mathbf{a}_{i}$). In view
of
\[
\prod_{i\in M}\mathbf{a}_{i}=\left(  \prod_{i\in U}\mathbf{a}_{i}\right)
\left(  \prod_{i\in V}\mathbf{a}_{i}\right)  \ \ \ \ \ \ \ \ \ \ \left(
\begin{array}
[c]{c}%
\text{since the set }M\text{ is the union of its}\\
\text{two disjoint subsets }U\text{ and }V
\end{array}
\right)  ,
\]
this rewrites as
\[
\left(  \prod_{i\in J}\mathbf{a}_{i}\right)  \left(  \prod_{i\in I\setminus
J}\mathbf{a}_{i}\right)  \overset{x^{n}}{\equiv}\prod_{i\in M}\mathbf{a}_{i}.
\]
In other words,
\[
\text{each }m\in\left\{  0,1,\ldots,n\right\}  \text{ satisfies }\left[
x^{m}\right]  \left(  \left(  \prod_{i\in J}\mathbf{a}_{i}\right)  \left(
\prod_{i\in I\setminus J}\mathbf{a}_{i}\right)  \right)  =\left[
x^{m}\right]  \left(  \prod_{i\in M}\mathbf{a}_{i}\right)
\]
(by Definition \ref{def.fps.xneq}). Applying this to $m=n$, we obtain
\[
\left[  x^{n}\right]  \left(  \left(  \prod_{i\in J}\mathbf{a}_{i}\right)
\left(  \prod_{i\in I\setminus J}\mathbf{a}_{i}\right)  \right)  =\left[
x^{n}\right]  \left(  \prod_{i\in M}\mathbf{a}_{i}\right)  .
\]
Comparing this with (\ref{pf.prop.fps.union-mulable.b.2}), we obtain%
\begin{equation}
\left[  x^{n}\right]  \left(  \prod_{i\in I}\mathbf{a}_{i}\right)  =\left[
x^{n}\right]  \left(  \left(  \prod_{i\in J}\mathbf{a}_{i}\right)  \left(
\prod_{i\in I\setminus J}\mathbf{a}_{i}\right)  \right)  .
\label{pf.prop.fps.union-mulable.b.at}%
\end{equation}


Now, forget that we fixed $n$. We thus have proved that each $n\in\mathbb{N}$
satisfies (\ref{pf.prop.fps.union-mulable.b.at}). In other words, each
coefficient of the FPS $\prod_{i\in I}\mathbf{a}_{i}$ equals the corresponding
coefficient of $\left(  \prod_{i\in J}\mathbf{a}_{i}\right)  \left(
\prod_{i\in I\setminus J}\mathbf{a}_{i}\right)  $. Hence, we have%
\[
\prod_{i\in I}\mathbf{a}_{i}=\left(  \prod_{i\in J}\mathbf{a}_{i}\right)
\cdot\left(  \prod_{i\in I\setminus J}\mathbf{a}_{i}\right)  .
\]
This proves Proposition \ref{prop.fps.union-mulable} \textbf{(b)}.
\end{proof}
\end{fineprint}

\begin{fineprint}
\begin{proof}
[Detailed proof of Proposition \ref{prop.fps.prod-mulable}.]\textbf{(a)} Fix
$n\in\mathbb{N}$. We know that the family $\left(  \mathbf{a}_{i}\right)
_{i\in I}$ is multipliable. Hence, there exists an $x^{n}$-approximator $U$
for $\left(  \mathbf{a}_{i}\right)  _{i\in I}$ (by Lemma
\ref{lem.fps.mulable.approx}). Consider this $U$.

We also know that the family $\left(  \mathbf{b}_{i}\right)  _{i\in I}$ is
multipliable. Hence, there exists an $x^{n}$-approximator $V$ for $\left(
\mathbf{b}_{i}\right)  _{i\in I}$ (by Lemma \ref{lem.fps.mulable.approx},
applied to $\mathbf{b}_{i}$ instead of $\mathbf{a}_{i}$). Consider this $V$.

We know that $U$ is an $x^{n}$-approximator for $\left(  \mathbf{a}%
_{i}\right)  _{i\in I}$. In other words, $U$ is a finite subset of $I$ that
determines the first $n+1$ coefficients in the product of $\left(
\mathbf{a}_{i}\right)  _{i\in I}$ (by the definition of an $x^{n}%
$-approximator). Hence, in particular, $U$ is finite. Similarly, $V$ is
finite. Moreover, $U\subseteq I$ (since $U$ is a subset of $I$); similarly,
$V\subseteq I$.

Let $M=U\cup V$. This set $M$ is finite (since $U$ and $V$ are finite).
Moreover, using $U\subseteq I$ and $V\subseteq I$, we obtain $M=\underbrace{U}%
_{\subseteq I}\cup\underbrace{V}_{\subseteq I}\subseteq I\cup I=I$. Hence, $M$
is a finite subset of $I$.

Now, let $N$ be a finite subset of $I$ satisfying $M\subseteq N\subseteq I$.
We shall show that
\[
\left[  x^{n}\right]  \left(  \prod_{i\in N}\left(  \mathbf{a}_{i}%
\mathbf{b}_{i}\right)  \right)  =\left[  x^{n}\right]  \left(  \prod_{i\in
M}\left(  \mathbf{a}_{i}\mathbf{b}_{i}\right)  \right)  .
\]


Indeed, let $m\in\left\{  0,1,\ldots,n\right\}  $. We have $U\subseteq U\cup
V=M\subseteq N$. The set $N$ is a finite subset of $I$ and satisfies
$U\subseteq N$. Now, recall that $U$ determines the first $n+1$ coefficients
in the product of $\left(  \mathbf{a}_{i}\right)  _{i\in I}$. Hence, $U$
determines the $x^{m}$-coefficient in the product of $\left(  \mathbf{a}%
_{i}\right)  _{i\in I}$ (since $m\in\left\{  0,1,\ldots,n\right\}  $). In
other words, every finite subset $T$ of $I$ satisfying $U\subseteq T\subseteq
I$ satisfies%
\begin{equation}
\left[  x^{m}\right]  \left(  \prod_{i\in T}\mathbf{a}_{i}\right)  =\left[
x^{m}\right]  \left(  \prod_{i\in U}\mathbf{a}_{i}\right)
\label{pf.prop.fps.prod-mulable.4}%
\end{equation}
(by the definition of what it means to \textquotedblleft determine the $x^{m}%
$-coefficient in the product of $\left(  \mathbf{a}_{i}\right)  _{i\in I}%
$\textquotedblright). We can apply this to $T=N$ (since $N$ is a finite subset
of $I$ satisfying $U\subseteq N\subseteq I$), and thus obtain%
\[
\left[  x^{m}\right]  \left(  \prod_{i\in N}\mathbf{a}_{i}\right)  =\left[
x^{m}\right]  \left(  \prod_{i\in U}\mathbf{a}_{i}\right)  .
\]
However, we can also apply (\ref{pf.prop.fps.prod-mulable.4}) to $T=M$ (since
$M$ is a finite subset of $I$ satisfying $U\subseteq M\subseteq I$), and thus
obtain%
\[
\left[  x^{m}\right]  \left(  \prod_{i\in M}\mathbf{a}_{i}\right)  =\left[
x^{m}\right]  \left(  \prod_{i\in U}\mathbf{a}_{i}\right)  .
\]
Comparing these two equalities, we obtain%
\begin{equation}
\left[  x^{m}\right]  \left(  \prod_{i\in N}\mathbf{a}_{i}\right)  =\left[
x^{m}\right]  \left(  \prod_{i\in M}\mathbf{a}_{i}\right)  .
\label{pf.prop.fps.prod-mulable.5a}%
\end{equation}
The same argument (applied to $\mathbf{b}_{i}$ and $V$ instead of
$\mathbf{a}_{i}$ and $U$) yields%
\begin{equation}
\left[  x^{m}\right]  \left(  \prod_{i\in N}\mathbf{b}_{i}\right)  =\left[
x^{m}\right]  \left(  \prod_{i\in M}\mathbf{b}_{i}\right)  .
\label{pf.prop.fps.prod-mulable.5b}%
\end{equation}


Forget that we fixed $m$. We thus have proved the equalities
(\ref{pf.prop.fps.prod-mulable.5a}) and (\ref{pf.prop.fps.prod-mulable.5b})
for each $m\in\left\{  0,1,\ldots,n\right\}  $. Hence, Lemma
\ref{lem.fps.prod.irlv.cong-mul} (applied to $a=\prod_{i\in N}\mathbf{a}_{i}$
and $b=\prod_{i\in M}\mathbf{a}_{i}$ and $c=\prod_{i\in N}\mathbf{b}_{i}$ and
$d=\prod_{i\in M}\mathbf{b}_{i}$) yields that%
\[
\left[  x^{m}\right]  \left(  \left(  \prod_{i\in N}\mathbf{a}_{i}\right)
\left(  \prod_{i\in N}\mathbf{b}_{i}\right)  \right)  =\left[  x^{m}\right]
\left(  \left(  \prod_{i\in M}\mathbf{a}_{i}\right)  \left(  \prod_{i\in
M}\mathbf{b}_{i}\right)  \right)
\]
for each $m\in\left\{  0,1,\ldots,n\right\}  $. Applying this to $m=n$, we
obtain%
\[
\left[  x^{n}\right]  \left(  \left(  \prod_{i\in N}\mathbf{a}_{i}\right)
\left(  \prod_{i\in N}\mathbf{b}_{i}\right)  \right)  =\left[  x^{n}\right]
\left(  \left(  \prod_{i\in M}\mathbf{a}_{i}\right)  \left(  \prod_{i\in
M}\mathbf{b}_{i}\right)  \right)  .
\]
In view of
\[
\prod_{i\in N}\left(  \mathbf{a}_{i}\mathbf{b}_{i}\right)  =\left(
\prod_{i\in N}\mathbf{a}_{i}\right)  \left(  \prod_{i\in N}\mathbf{b}%
_{i}\right)  \ \ \ \ \ \ \ \ \ \ \left(
\begin{array}
[c]{c}%
\text{by the standard rules for finite}\\
\text{products, since }N\text{ is a finite set}%
\end{array}
\right)
\]
and%
\[
\prod_{i\in M}\left(  \mathbf{a}_{i}\mathbf{b}_{i}\right)  =\left(
\prod_{i\in M}\mathbf{a}_{i}\right)  \left(  \prod_{i\in M}\mathbf{b}%
_{i}\right)  \ \ \ \ \ \ \ \ \ \ \left(
\begin{array}
[c]{c}%
\text{by the standard rules for finite}\\
\text{products, since }M\text{ is a finite set}%
\end{array}
\right)  ,
\]
this rewrites as%
\[
\left[  x^{n}\right]  \left(  \prod_{i\in N}\left(  \mathbf{a}_{i}%
\mathbf{b}_{i}\right)  \right)  =\left[  x^{n}\right]  \left(  \prod_{i\in
M}\left(  \mathbf{a}_{i}\mathbf{b}_{i}\right)  \right)  .
\]


Forget that we fixed $N$. We thus have shown that every finite subset $N$ of
$I$ satisfying $M\subseteq N\subseteq I$ satisfies $\left[  x^{n}\right]
\left(  \prod_{i\in N}\left(  \mathbf{a}_{i}\mathbf{b}_{i}\right)  \right)
=\left[  x^{n}\right]  \left(  \prod_{i\in M}\left(  \mathbf{a}_{i}%
\mathbf{b}_{i}\right)  \right)  $. In other words, $M$ determines the $x^{n}%
$-coefficient in the product of $\left(  \mathbf{a}_{i}\mathbf{b}_{i}\right)
_{i\in I}$ (by the definition of what it means to \textquotedblleft determine
the $x^{n}$-coefficient in the product of $\left(  \mathbf{a}_{i}%
\mathbf{b}_{i}\right)  _{i\in I}$\textquotedblright). Hence, the $x^{n}%
$-coefficient in the product of $\left(  \mathbf{a}_{i}\mathbf{b}_{i}\right)
_{i\in I}$ is finitely determined (by the definition of \textquotedblleft
finitely determined\textquotedblright, since $M$ is a finite subset of $I$).

Forget that we fixed $n$. We thus have proved that for each $n\in\mathbb{N}$,
the $x^{n}$-coefficient in the product of $\left(  \mathbf{a}_{i}%
\mathbf{b}_{i}\right)  _{i\in I}$ is finitely determined. In other words, each
coefficient in the product of $\left(  \mathbf{a}_{i}\mathbf{b}_{i}\right)
_{i\in I}$ is finitely determined. In other words, the family $\left(
\mathbf{a}_{i}\mathbf{b}_{i}\right)  _{i\in I}$ is multipliable (by the
definition of \textquotedblleft multipliable\textquotedblright). This proves
Proposition \ref{prop.fps.prod-mulable} \textbf{(a)}. \medskip

\textbf{(b)} Proposition \ref{prop.fps.prod-mulable} \textbf{(a)} shows that
the family $\left(  \mathbf{a}_{i}\mathbf{b}_{i}\right)  _{i\in I}$ is multipliable.

Let $n\in\mathbb{N}$. In our above proof of Proposition
\ref{prop.fps.prod-mulable} \textbf{(a)}, we have seen the following:

\begin{itemize}
\item There exists an $x^{n}$-approximator $U$ for $\left(  \mathbf{a}%
_{i}\right)  _{i\in I}$.

\item There exists an $x^{n}$-approximator $V$ for $\left(  \mathbf{b}%
_{i}\right)  _{i\in I}$.
\end{itemize}

Consider these $U$ and $V$. Let $M=U\cup V$. Thus, $M=U\cup V\supseteq U$ and
$M=U\cup V\supseteq V$. In our above proof of Proposition
\ref{prop.fps.prod-mulable} \textbf{(a)}, we have seen the following:

\begin{itemize}
\item The set $M$ is a finite subset of $I$.

\item The set $M$ determines the $x^{n}$-coefficient in the product of
$\left(  \mathbf{a}_{i}\mathbf{b}_{i}\right)  _{i\in I}$.
\end{itemize}

We recall that the relation $\overset{x^{n}}{\equiv}$ on $K\left[  \left[
x\right]  \right]  $ is an equivalence relation (by Theorem
\ref{thm.fps.xneq.props} \textbf{(a)}). Thus, this relation $\overset{x^{n}%
}{\equiv}$ is transitive and symmetric.

Now, the definition of the infinite product $\prod_{i\in I}\left(
\mathbf{a}_{i}\mathbf{b}_{i}\right)  $ (namely, Definition
\ref{def.fps.multipliable} \textbf{(b)}) yields that%
\begin{equation}
\left[  x^{n}\right]  \left(  \prod_{i\in I}\left(  \mathbf{a}_{i}%
\mathbf{b}_{i}\right)  \right)  =\left[  x^{n}\right]  \left(  \prod_{i\in
M}\left(  \mathbf{a}_{i}\mathbf{b}_{i}\right)  \right)
\label{pf.prop.fps.prod-mulable.b.2}%
\end{equation}
(since $M$ is a finite subset of $I$ that determines the $x^{n}$-coefficient
in the product of $\left(  \mathbf{a}_{i}\mathbf{b}_{i}\right)  _{i\in I}$).
On the other hand, the set $U$ is an $x^{n}$-approximator for $\left(
\mathbf{a}_{i}\right)  _{i\in I}$. Thus, Proposition
\ref{prop.fps.infprod-approx-xneq} \textbf{(b)} (applied to $U$ instead of
$M$) yields
\[
\prod_{i\in I}\mathbf{a}_{i}\overset{x^{n}}{\equiv}\prod_{i\in U}%
\mathbf{a}_{i}%
\]
(since the family $\left(  \mathbf{a}_{i}\right)  _{i\in I}$ is multipliable).
Since the relation $\overset{x^{n}}{\equiv}$ is symmetric, we thus obtain%
\begin{equation}
\prod_{i\in U}\mathbf{a}_{i}\overset{x^{n}}{\equiv}\prod_{i\in I}%
\mathbf{a}_{i}. \label{pf.prop.fps.prod-mulable.b.3}%
\end{equation}
Moreover, $M$ is a finite subset of $I$ satisfying $U\subseteq M$ (since
$M\supseteq U$) and therefore $U\subseteq M\subseteq I$; hence, Proposition
\ref{prop.fps.infprod-approx-xneq} \textbf{(a)} (applied to $U$ and $M$
instead of $M$ and $J$) yields%
\[
\prod_{i\in M}\mathbf{a}_{i}\overset{x^{n}}{\equiv}\prod_{i\in U}%
\mathbf{a}_{i}\overset{x^{n}}{\equiv}\prod_{i\in I}\mathbf{a}_{i}%
\ \ \ \ \ \ \ \ \ \ \left(  \text{by (\ref{pf.prop.fps.prod-mulable.b.3}%
)}\right)  .
\]
Hence,%
\begin{equation}
\prod_{i\in M}\mathbf{a}_{i}\overset{x^{n}}{\equiv}\prod_{i\in I}%
\mathbf{a}_{i} \label{pf.prop.fps.prod-mulable.b.4}%
\end{equation}
(since the relation $\overset{x^{n}}{\equiv}$ is transitive). The same
argument (applied to $\left(  \mathbf{b}_{i}\right)  _{i\in I}$ and $V$
instead of $\left(  \mathbf{a}_{i}\right)  _{i\in I}$ and $U$) yields%
\begin{equation}
\prod_{i\in M}\mathbf{b}_{i}\overset{x^{n}}{\equiv}\prod_{i\in I}%
\mathbf{b}_{i}. \label{pf.prop.fps.prod-mulable.b.5}%
\end{equation}
From (\ref{pf.prop.fps.prod-mulable.b.4}) and
(\ref{pf.prop.fps.prod-mulable.b.5}), we obtain%
\[
\left(  \prod_{i\in M}\mathbf{a}_{i}\right)  \left(  \prod_{i\in M}%
\mathbf{b}_{i}\right)  \overset{x^{n}}{\equiv}\left(  \prod_{i\in I}%
\mathbf{a}_{i}\right)  \left(  \prod_{i\in I}\mathbf{b}_{i}\right)
\]
(by (\ref{eq.thm.fps.xneq.props.b.*}), applied to $a=\prod_{i\in M}%
\mathbf{a}_{i}$ and $b=\prod_{i\in I}\mathbf{a}_{i}$ and $c=\prod_{i\in
M}\mathbf{b}_{i}$ and $d=\prod_{i\in I}\mathbf{b}_{i}$). In view of
\[
\left(  \prod_{i\in M}\mathbf{a}_{i}\right)  \left(  \prod_{i\in M}%
\mathbf{b}_{i}\right)  =\prod_{i\in M}\left(  \mathbf{a}_{i}\mathbf{b}%
_{i}\right)  \ \ \ \ \ \ \ \ \ \ \left(
\begin{array}
[c]{c}%
\text{by the properties of finite products,}\\
\text{since the set }M\text{ is finite}%
\end{array}
\right)  ,
\]
this rewrites as
\[
\prod_{i\in M}\left(  \mathbf{a}_{i}\mathbf{b}_{i}\right)  \overset{x^{n}%
}{\equiv}\left(  \prod_{i\in I}\mathbf{a}_{i}\right)  \left(  \prod_{i\in
I}\mathbf{b}_{i}\right)  .
\]
In other words,
\[
\text{each }m\in\left\{  0,1,\ldots,n\right\}  \text{ satisfies }\left[
x^{m}\right]  \left(  \prod_{i\in M}\left(  \mathbf{a}_{i}\mathbf{b}%
_{i}\right)  \right)  =\left[  x^{m}\right]  \left(  \left(  \prod_{i\in
I}\mathbf{a}_{i}\right)  \left(  \prod_{i\in I}\mathbf{b}_{i}\right)  \right)
\]
(by Definition \ref{def.fps.xneq}). Applying this to $m=n$, we obtain
\[
\left[  x^{n}\right]  \left(  \prod_{i\in M}\left(  \mathbf{a}_{i}%
\mathbf{b}_{i}\right)  \right)  =\left[  x^{n}\right]  \left(  \left(
\prod_{i\in I}\mathbf{a}_{i}\right)  \left(  \prod_{i\in I}\mathbf{b}%
_{i}\right)  \right)  .
\]
In view of (\ref{pf.prop.fps.prod-mulable.b.2}), this rewrites as%
\begin{equation}
\left[  x^{n}\right]  \left(  \prod_{i\in I}\left(  \mathbf{a}_{i}%
\mathbf{b}_{i}\right)  \right)  =\left[  x^{n}\right]  \left(  \left(
\prod_{i\in I}\mathbf{a}_{i}\right)  \left(  \prod_{i\in I}\mathbf{b}%
_{i}\right)  \right)  . \label{pf.prop.fps.prod-mulable.b.at}%
\end{equation}


Now, forget that we fixed $n$. We thus have proved that each $n\in\mathbb{N}$
satisfies (\ref{pf.prop.fps.prod-mulable.b.at}). In other words, each
coefficient of the FPS $\prod_{i\in I}\left(  \mathbf{a}_{i}\mathbf{b}%
_{i}\right)  $ equals the corresponding coefficient of $\left(  \prod_{i\in
I}\mathbf{a}_{i}\right)  \left(  \prod_{i\in I}\mathbf{b}_{i}\right)  $.
Hence, we have%
\[
\prod_{i\in I}\left(  \mathbf{a}_{i}\mathbf{b}_{i}\right)  =\left(
\prod_{i\in I}\mathbf{a}_{i}\right)  \left(  \prod_{i\in I}\mathbf{b}%
_{i}\right)  .
\]
This proves Proposition \ref{prop.fps.prod-mulable} \textbf{(b)}.
\end{proof}
\end{fineprint}

\begin{fineprint}
In order to prove Proposition \ref{prop.fps.div-mulable}, we need a lemma (an
analogue of Lemma \ref{lem.fps.prod.irlv.cong-mul} for division instead of multiplication):

\begin{lemma}
\label{lem.fps.prod.irlv.cong-div}Let $a,b,c,d\in K\left[  \left[  x\right]
\right]  $ be four FPSs such that $c$ and $d$ are invertible. Let
$n\in\mathbb{N}$. Assume that
\[
\left[  x^{m}\right]  a=\left[  x^{m}\right]  b\ \ \ \ \ \ \ \ \ \ \text{for
each }m\in\left\{  0,1,\ldots,n\right\}  .
\]
Assume further that%
\[
\left[  x^{m}\right]  c=\left[  x^{m}\right]  d\ \ \ \ \ \ \ \ \ \ \text{for
each }m\in\left\{  0,1,\ldots,n\right\}  .
\]
Then,
\[
\left[  x^{m}\right]  \dfrac{a}{c}=\left[  x^{m}\right]  \dfrac{b}%
{d}\ \ \ \ \ \ \ \ \ \ \text{for each }m\in\left\{  0,1,\ldots,n\right\}  .
\]

\end{lemma}

\begin{proof}
[Proof of Lemma \ref{lem.fps.prod.irlv.cong-div}.]We have assumed that
\[
\left[  x^{m}\right]  a=\left[  x^{m}\right]  b\ \ \ \ \ \ \ \ \ \ \text{for
each }m\in\left\{  0,1,\ldots,n\right\}  .
\]
In other words, $a\overset{x^{n}}{\equiv}b$ (by the definition of the relation
$\overset{x^{n}}{\equiv}$). Similarly, $c\overset{x^{n}}{\equiv}d$. Hence,
(\ref{eq.thm.fps.xneq.props.e./}) yields $\dfrac{a}{c}\overset{x^{n}}{\equiv
}\dfrac{b}{d}$. In other words,%
\[
\left[  x^{m}\right]  \dfrac{a}{c}=\left[  x^{m}\right]  \dfrac{b}%
{d}\ \ \ \ \ \ \ \ \ \ \text{for each }m\in\left\{  0,1,\ldots,n\right\}
\]
(by the definition of the relation $\overset{x^{n}}{\equiv}$). This proves
Lemma \ref{lem.fps.prod.irlv.cong-div}.
\end{proof}

Now we can easily prove Proposition \ref{prop.fps.div-mulable}:

\begin{proof}
[Detailed proof of Proposition \ref{prop.fps.div-mulable}.]This can be proved
similarly to Proposition \ref{prop.fps.prod-mulable}, with the obvious changes
(replacing multiplication by division, and requiring each $\mathbf{b}_{i}$ to
be invertible). Of course, instead of Lemma \ref{lem.fps.prod.irlv.cong-mul},
we need to use Lemma \ref{lem.fps.prod.irlv.cong-div}.
\end{proof}
\end{fineprint}

\begin{fineprint}
In order to eventually prove Proposition \ref{prop.fps.prods-mulable-subfams},
we shall first prove a slightly stronger auxiliary statement:

\begin{lemma}
\label{lem.fps.prods-mulable-subfams-appr}Let $\left(  \mathbf{a}_{i}\right)
_{i\in I}\in K\left[  \left[  x\right]  \right]  ^{I}$ be a family of
invertible FPSs. Let $J$ be a subset of $I$. Let $n\in\mathbb{N}$. Let $U$ be
an $x^{n}$-approximator for $\left(  \mathbf{a}_{i}\right)  _{i\in I}$. Then,
$U\cap J$ is an $x^{n}$-approximator for $\left(  \mathbf{a}_{i}\right)
_{i\in J}$.
\end{lemma}

\begin{proof}
[Proof of Lemma \ref{lem.fps.prods-mulable-subfams-appr}.]The set $U$ is an
$x^{n}$-approximator for $\left(  \mathbf{a}_{i}\right)  _{i\in I}$. In other
words, $U$ is a finite subset of $I$ that determines the first $n+1$
coefficients in the product of $\left(  \mathbf{a}_{i}\right)  _{i\in I}$ (by
the definition of an $x^{n}$-approximator). Hence, in particular, $U$ is
finite. Moreover, $U\subseteq I$ (since $U$ is a subset of $I$).

Let $M=U\cap J$. Thus, $M=U\cap J\subseteq U$, so that the set $M$ is finite
(since $U$ is finite). Moreover, $M$ is a subset of $J$ (since $M=U\cap
J\subseteq J$).

Now, let $N$ be a finite subset of $J$ satisfying $M\subseteq N\subseteq J$.
We shall show that
\[
\prod_{i\in N}\mathbf{a}_{i}\overset{x^{n}}{\equiv}\prod_{i\in M}%
\mathbf{a}_{i}.
\]


Indeed, the set $N\cup U$ is finite (since $N$ and $U$ are finite) and is a
subset of $I$ (since $\underbrace{N}_{\subseteq J\subseteq I}\cup
\underbrace{U}_{\subseteq I}\subseteq I\cup I=I$); it also satisfies
$U\subseteq N\cup U\subseteq I$. Now, we recall that $U$ is an $x^{n}%
$-approximator for $\left(  \mathbf{a}_{i}\right)  _{i\in I}$. Hence,
Proposition \ref{prop.fps.infprod-approx-xneq} \textbf{(a)} (applied to $U$
and $N\cup U$ instead of $M$ and $J$) yields
\begin{equation}
\prod_{i\in N\cup U}\mathbf{a}_{i}\overset{x^{n}}{\equiv}\prod_{i\in
U}\mathbf{a}_{i} \label{pf.lem.fps.prods-mulable-subfams-appr.ab}%
\end{equation}
(since $N\cup U$ is a finite subset of $I$ satisfying $U\subseteq N\cup
U\subseteq I$).

On the other hand, we have assumed that $\left(  \mathbf{a}_{i}\right)  _{i\in
I}$ is a family of invertible FPSs. Thus, for each $i\in I$, the FPS
$\mathbf{a}_{i}$ is invertible, so that its inverse $\mathbf{a}_{i}^{-1}$ is
well-defined. We obviously have%
\begin{equation}
\prod_{i\in U\setminus J}\mathbf{a}_{i}^{-1}\overset{x^{n}}{\equiv}\prod_{i\in
U\setminus J}\mathbf{a}_{i}^{-1}
\label{pf.lem.fps.prods-mulable-subfams-appr.cd}%
\end{equation}
(since Proposition \ref{thm.fps.xneq.props} \textbf{(a)} yields that the
relation $\overset{x^{n}}{\equiv}$ is reflexive).

We have now proved (\ref{pf.lem.fps.prods-mulable-subfams-appr.ab}) and
(\ref{pf.lem.fps.prods-mulable-subfams-appr.cd}). Hence,
(\ref{eq.thm.fps.xneq.props.b.*}) (applied to $a=\prod_{i\in N\cup
U}\mathbf{a}_{i}$ and $b=\prod_{i\in U}\mathbf{a}_{i}$ and $c=\prod_{i\in
U\setminus J}\mathbf{a}_{i}^{-1}$ and $d=\prod_{i\in U\setminus J}%
\mathbf{a}_{i}^{-1}$) yields%
\begin{equation}
\left(  \prod_{i\in N\cup U}\mathbf{a}_{i}\right)  \left(  \prod_{i\in
U\setminus J}\mathbf{a}_{i}^{-1}\right)  \overset{x^{n}}{\equiv}\left(
\prod_{i\in U}\mathbf{a}_{i}\right)  \left(  \prod_{i\in U\setminus
J}\mathbf{a}_{i}^{-1}\right)  .
\label{pf.lem.fps.prods-mulable-subfams-appr.ac=bd}%
\end{equation}


On the other hand, the sets $N$ and $U\setminus J$ are disjoint\footnote{since
$\underbrace{N}_{\subseteq J}\cap\left(  U\setminus J\right)  \subseteq
J\cap\left(  U\setminus J\right)  =\varnothing$ and thus $N\cap\left(
U\setminus J\right)  =\varnothing$}, and their union is $N\cup\left(
U\setminus J\right)  =N\cup U$\ \ \ \ \footnote{This follows from%
\[
N\cup\underbrace{U}_{=\left(  U\cap J\right)  \cup\left(  U\setminus J\right)
}=N\cup\underbrace{\left(  U\cap J\right)  }_{=M}\cup\left(  U\setminus
J\right)  \subseteq\underbrace{N\cup M}_{\substack{=N\\\text{(since
}M\subseteq N\text{)}}}\cup\left(  U\setminus J\right)  =N\cup\left(
U\setminus J\right)  .
\]
}. Hence, the set $N\cup U$ is the union of its two disjoint subsets $N$ and
$U\setminus J$. Thus, we can split the product $\prod_{i\in N\cup U}%
\mathbf{a}_{i}$ as follows:%
\[
\prod_{i\in N\cup U}\mathbf{a}_{i}=\left(  \prod_{i\in N}\mathbf{a}%
_{i}\right)  \left(  \prod_{i\in U\setminus J}\mathbf{a}_{i}\right)  .
\]
Multiplying both sides of this equality by $\prod_{i\in U\setminus
J}\mathbf{a}_{i}^{-1}$, we obtain%
\begin{align}
\left(  \prod_{i\in N\cup U}\mathbf{a}_{i}\right)  \left(  \prod_{i\in
U\setminus J}\mathbf{a}_{i}^{-1}\right)   &  =\left(  \prod_{i\in N}%
\mathbf{a}_{i}\right)  \underbrace{\left(  \prod_{i\in U\setminus J}%
\mathbf{a}_{i}\right)  \left(  \prod_{i\in U\setminus J}\mathbf{a}_{i}%
^{-1}\right)  }_{=\prod_{i\in U\setminus J}\left(  \mathbf{a}_{i}%
\mathbf{a}_{i}^{-1}\right)  }\nonumber\\
&  =\left(  \prod_{i\in N}\mathbf{a}_{i}\right)  \left(  \prod_{i\in
U\setminus J}\underbrace{\left(  \mathbf{a}_{i}\mathbf{a}_{i}^{-1}\right)
}_{=1}\right)  =\left(  \prod_{i\in N}\mathbf{a}_{i}\right)
\underbrace{\left(  \prod_{i\in U\setminus J}1\right)  }_{=1}\nonumber\\
&  =\prod_{i\in N}\mathbf{a}_{i}. \label{pf.prop.fps.prods-mulable-subfams.7a}%
\end{align}


However, the set $U$ is the union of its two disjoint subsets $U\cap J$ and
$U\setminus J$. Thus, we can split the product $\prod_{i\in U}\mathbf{a}_{i}$
as follows:%
\[
\prod_{i\in U}\mathbf{a}_{i}=\left(  \prod_{i\in U\cap J}\mathbf{a}%
_{i}\right)  \left(  \prod_{i\in U\setminus J}\mathbf{a}_{i}\right)  .
\]
Multiplying both sides of this equality by $\prod_{i\in U\setminus
J}\mathbf{a}_{i}^{-1}$, we obtain%
\begin{align}
\left(  \prod_{i\in U}\mathbf{a}_{i}\right)  \left(  \prod_{i\in U\setminus
J}\mathbf{a}_{i}^{-1}\right)   &  =\underbrace{\left(  \prod_{i\in U\cap
J}\mathbf{a}_{i}\right)  }_{\substack{=\prod_{i\in M}\mathbf{a}_{i}%
\\\text{(since }U\cap J=M\text{)}}}\underbrace{\left(  \prod_{i\in U\setminus
J}\mathbf{a}_{i}\right)  \left(  \prod_{i\in U\setminus J}\mathbf{a}_{i}%
^{-1}\right)  }_{=\prod_{i\in U\setminus J}\left(  \mathbf{a}_{i}%
\mathbf{a}_{i}^{-1}\right)  }\nonumber\\
&  =\left(  \prod_{i\in M}\mathbf{a}_{i}\right)  \left(  \prod_{i\in
U\setminus J}\underbrace{\left(  \mathbf{a}_{i}\mathbf{a}_{i}^{-1}\right)
}_{=1}\right)  =\left(  \prod_{i\in M}\mathbf{a}_{i}\right)
\underbrace{\left(  \prod_{i\in U\setminus J}1\right)  }_{=1}\nonumber\\
&  =\prod_{i\in M}\mathbf{a}_{i}. \label{pf.prop.fps.prods-mulable-subfams.7b}%
\end{align}


In view of (\ref{pf.prop.fps.prods-mulable-subfams.7a}) and
(\ref{pf.prop.fps.prods-mulable-subfams.7b}), we can rewrite the relation
(\ref{pf.lem.fps.prods-mulable-subfams-appr.ac=bd}) as follows:%
\[
\prod_{i\in N}\mathbf{a}_{i}\overset{x^{n}}{\equiv}\prod_{i\in M}%
\mathbf{a}_{i}.
\]
In other words, each $m\in\left\{  0,1,\ldots,n\right\}  $ satisfies
\begin{equation}
\left[  x^{m}\right]  \left(  \prod_{i\in N}\mathbf{a}_{i}\right)  =\left[
x^{m}\right]  \left(  \prod_{i\in M}\mathbf{a}_{i}\right)
\label{pf.lem.fps.prods-mulable-subfams-appr.at}%
\end{equation}
(by Definition \ref{def.fps.xneq}).

Forget that we fixed $N$. We have thus shown that every finite subset $N$ of
$J$ satisfying $M\subseteq N\subseteq J$ satisfies
(\ref{pf.lem.fps.prods-mulable-subfams-appr.at}) for each $m\in\left\{
0,1,\ldots,n\right\}  $.

Now, let $m\in\left\{  0,1,\ldots,n\right\}  $. Then, every finite subset $N$
of $J$ satisfying $M\subseteq N\subseteq J$ satisfies%
\[
\left[  x^{m}\right]  \left(  \prod_{i\in N}\mathbf{a}_{i}\right)  =\left[
x^{m}\right]  \left(  \prod_{i\in M}\mathbf{a}_{i}\right)
\]
(by (\ref{pf.lem.fps.prods-mulable-subfams-appr.at})). In other words, the set
$M$ determines the $x^{m}$-coefficient in the product of $\left(
\mathbf{a}_{i}\right)  _{i\in J}$ (by the definition of \textquotedblleft
determining the $x^{m}$-coefficient in a product\textquotedblright).

Forget that we fixed $m$. We thus have shown that $M$ determines the $x^{m}%
$-coefficient in the product of $\left(  \mathbf{a}_{i}\right)  _{i\in J}$ for
each $m\in\left\{  0,1,\ldots,n\right\}  $. In other words, $M$ determines the
first $n+1$ coefficients in the product of $\left(  \mathbf{a}_{i}\right)
_{i\in J}$. In other words, $M$ is an $x^{n}$-approximator for $\left(
\mathbf{a}_{i}\right)  _{i\in J}$ (by the definition of an \textquotedblleft%
$x^{n}$-approximator\textquotedblright, since $M$ is a finite subset of $J$).
In other words, $U\cap J$ is an $x^{n}$-approximator for $\left(
\mathbf{a}_{i}\right)  _{i\in J}$ (since $M=U\cap J$). This proves Lemma
\ref{lem.fps.prods-mulable-subfams-appr}.
\end{proof}
\end{fineprint}

\begin{fineprint}
\begin{proof}
[Detailed proof of Proposition \ref{prop.fps.prods-mulable-subfams}.]Let $J$
be a subset of $I$. We shall show that the family $\left(  \mathbf{a}%
_{i}\right)  _{i\in J}$ is multipliable.

Fix $n\in\mathbb{N}$. We know that the family $\left(  \mathbf{a}_{i}\right)
_{i\in I}$ is multipliable. Hence, there exists an $x^{n}$-approximator $U$
for $\left(  \mathbf{a}_{i}\right)  _{i\in I}$ (by Lemma
\ref{lem.fps.mulable.approx}). Consider this $U$.

Set $M=U\cap J$. Lemma \ref{lem.fps.prods-mulable-subfams-appr} yields that
$U\cap J$ is an $x^{n}$-approximator for $\left(  \mathbf{a}_{i}\right)
_{i\in J}$. In other words, $M$ is an $x^{n}$-approximator for $\left(
\mathbf{a}_{i}\right)  _{i\in J}$ (since $M=U\cap J$). In other words, $M$ is
a finite subset of $J$ that determines the first $n+1$ coefficients in the
product of $\left(  \mathbf{a}_{i}\right)  _{i\in J}$ (by the definition of an
$x^{n}$-approximator). Thus, the set $M$ determines the first $n+1$
coefficients in the product of $\left(  \mathbf{a}_{i}\right)  _{i\in J}$.
Hence, in particular, this set $M$ determines the $x^{n}$-coefficient in the
product of $\left(  \mathbf{a}_{i}\right)  _{i\in J}$. Therefore, the $x^{n}%
$-coefficient in the product of $\left(  \mathbf{a}_{i}\right)  _{i\in J}$ is
finitely determined (by the definition of \textquotedblleft finitely
determined\textquotedblright, since $M$ is a finite subset of $J$).

Forget that we fixed $n$. We thus have proved that for each $n\in\mathbb{N}$,
the $x^{n}$-coefficient in the product of $\left(  \mathbf{a}_{i}\right)
_{i\in J}$ is finitely determined. In other words, each coefficient in the
product of $\left(  \mathbf{a}_{i}\right)  _{i\in J}$ is finitely determined.
In other words, the family $\left(  \mathbf{a}_{i}\right)  _{i\in J}$ is
multipliable (by the definition of \textquotedblleft
multipliable\textquotedblright).

Forget that we fixed $J$. We thus have shown that the family $\left(
\mathbf{a}_{i}\right)  _{i\in J}$ is multipliable whenever $J$ is a subset of
$I$. In other words, any subfamily of $\left(  \mathbf{a}_{i}\right)  _{i\in
I}$ is multipliable. This proves Proposition
\ref{prop.fps.prods-mulable-subfams}.
\end{proof}
\end{fineprint}

\begin{fineprint}
Our next goal is to prove Proposition \ref{prop.fps.prods-mulable-rules.SW1}.
First, however, let us restate a piece of Theorem \ref{thm.fps.xneq.props}
\textbf{(f)} in more convenient language:

\begin{lemma}
\label{lem.fps.prods-mulable-rules.SW1.lem1}Let $n\in\mathbb{N}$. Let $V$ be a
finite set. Let $\left(  c_{w}\right)  _{w\in V}\in K\left[  \left[  x\right]
\right]  ^{V}$ and $\left(  d_{w}\right)  _{w\in V}\in K\left[  \left[
x\right]  \right]  ^{V}$ be two families of FPSs such that
\[
\text{each }w\in V\text{ satisfies }c_{w}\overset{x^{n}}{\equiv}d_{w}.
\]
Then, we have%
\[
\prod_{w\in V}c_{w}\overset{x^{n}}{\equiv}\prod_{w\in V}d_{w}.
\]

\end{lemma}

\begin{proof}
[Proof of Lemma \ref{lem.fps.prods-mulable-rules.SW1.lem1}.]This is just
(\ref{eq.thm.fps.xneq.props.e.*}), with the letters $S$, $s$, $a_{s}$ and
$b_{s}$ renamed as $V$, $w$, $c_{w}$ and $d_{w}$.
\end{proof}

\begin{proof}
[Detailed proof of Proposition \ref{prop.fps.prods-mulable-rules.SW1}.]We
shall subdivide our proof into several claims:

\begin{statement}
\textit{Claim 1:} Let $w\in W$. Then, the family $\left(  \mathbf{a}%
_{s}\right)  _{s\in S;\ f\left(  s\right)  =w}$ is multipliable.
\end{statement}

[\textit{Proof of Claim 1:} This was an assumption of Proposition
\ref{prop.fps.prods-mulable-rules.SW1}.] \medskip

Let us set%
\begin{equation}
\mathbf{b}_{w}:=\prod_{\substack{s\in S;\\f\left(  s\right)  =w}%
}\mathbf{a}_{s}\ \ \ \ \ \ \ \ \ \ \text{for each }w\in W.
\label{pf.prop.fps.prods-mulable-rules.SW1.bw=}%
\end{equation}
This is well-defined, because for each $w\in W$, the product $\prod
_{\substack{s\in S;\\f\left(  s\right)  =w}}\mathbf{a}_{s}$ is well-defined
(since Claim 1 shows that the family $\left(  \mathbf{a}_{s}\right)  _{s\in
S;\ f\left(  s\right)  =w}$ is multipliable).

Now, let $n\in\mathbb{N}$. Lemma \ref{lem.fps.mulable.approx} (applied to $S$
and $\left(  \mathbf{a}_{s}\right)  _{s\in S}$ instead of $I$ and $\left(
\mathbf{a}_{i}\right)  _{i\in I}$) shows that there exists an $x^{n}%
$-approximator for $\left(  \mathbf{a}_{s}\right)  _{s\in S}$. Pick such an
$x^{n}$-approximator, and call it $U$. Then, $U$ is an $x^{n}$-approximator
for $\left(  \mathbf{a}_{s}\right)  _{s\in S}$; in other words, $U$ is a
finite subset of $S$ that determines the first $n+1$ coefficients in the
product of $\left(  \mathbf{a}_{s}\right)  _{s\in S}$ (by the definition of an
$x^{n}$-approximator).

The set $U$ is finite. Thus, its image $f\left(  U\right)  =\left\{  f\left(
u\right)  \ \mid\ u\in U\right\}  $ is finite as well. Now, we claim the following:

\begin{statement}
\textit{Claim 2:} For each $w\in W$, we have%
\[
\mathbf{b}_{w}\overset{x^{n}}{\equiv}\prod_{\substack{s\in U;\\f\left(
s\right)  =w}}\mathbf{a}_{s}.
\]

\end{statement}

[\textit{Proof of Claim 2:} Let $w\in W$. Let $J$ be the subset $\left\{  s\in
S\ \mid\ f\left(  s\right)  =w\right\}  $ of $S$. Then, the family $\left(
\mathbf{a}_{s}\right)  _{s\in J}$ is just the family $\left(  \mathbf{a}%
_{s}\right)  _{s\in S;\ f\left(  s\right)  =w}$, and thus is multipliable (by
Claim 1). Furthermore, Lemma \ref{lem.fps.prods-mulable-subfams-appr} (applied
to $S$ and $\left(  \mathbf{a}_{s}\right)  _{s\in S}$ instead of $I$ and
$\left(  \mathbf{a}_{i}\right)  _{i\in I}$) yields that $U\cap J$ is an
$x^{n}$-approximator for $\left(  \mathbf{a}_{s}\right)  _{s\in J}$ (since $U$
is an $x^{n}$-approximator for $\left(  \mathbf{a}_{s}\right)  _{s\in S}$).
Hence, Proposition \ref{prop.fps.infprod-approx-xneq} \textbf{(b)} (applied to
$J$, $\left(  \mathbf{a}_{s}\right)  _{s\in J}$ and $U\cap J$ instead of $I$,
$\left(  \mathbf{a}_{i}\right)  _{i\in I}$ and $M$) yields
\[
\prod_{i\in J}\mathbf{a}_{i}\overset{x^{n}}{\equiv}\prod_{i\in U\cap
J}\mathbf{a}_{i}.
\]
Renaming the indices $i$ as $s$ on both sides of this relation, we obtain%
\begin{equation}
\prod_{s\in J}\mathbf{a}_{s}\overset{x^{n}}{\equiv}\prod_{s\in U\cap
J}\mathbf{a}_{s}. \label{pf.prop.fps.prods-mulable-rules.SW1.c2.pf.1}%
\end{equation}


However, we have $J=\left\{  s\in S\ \mid\ f\left(  s\right)  =w\right\}  $.
Thus, the product sign \textquotedblleft$\prod_{s\in J}$\textquotedblright\ is
equivalent to \textquotedblleft$\prod_{\substack{s\in S;\\f\left(  s\right)
=w}}$\textquotedblright. Thus, we obtain%
\begin{equation}
\prod_{s\in J}\mathbf{a}_{s}=\prod_{\substack{s\in S;\\f\left(  s\right)
=w}}\mathbf{a}_{s}=\mathbf{b}_{w}
\label{pf.prop.fps.prods-mulable-rules.SW1.c2.pf.2}%
\end{equation}
(by (\ref{pf.prop.fps.prods-mulable-rules.SW1.bw=})).

On the other hand, from $J=\left\{  s\in S\ \mid\ f\left(  s\right)
=w\right\}  $, we obtain%
\begin{align*}
U\cap J  &  =U\cap\left\{  s\in S\ \mid\ f\left(  s\right)  =w\right\} \\
&  =\left\{  s\in U\ \mid\ f\left(  s\right)  =w\right\}
\ \ \ \ \ \ \ \ \ \ \left(  \text{since }U\subseteq S\right)  .
\end{align*}
Hence, the product sign \textquotedblleft$\prod_{s\in U\cap J}$%
\textquotedblright\ is equivalent to \textquotedblleft$\prod_{\substack{s\in
U;\\f\left(  s\right)  =w}}$\textquotedblright. Thus, we obtain%
\begin{equation}
\prod_{s\in U\cap J}\mathbf{a}_{s}=\prod_{\substack{s\in U;\\f\left(
s\right)  =w}}\mathbf{a}_{s}.
\label{pf.prop.fps.prods-mulable-rules.SW1.c2.pf.3}%
\end{equation}